% ======================= 2. ASSUMPTIONS AND NOTATION =======================
\section{Assumptions and Notation}\label{sec:assume}

Every simplification below is stated with its rationale and with the direction in
which it pushes the answer, so that a reader can discount the results accordingly.
The AI workflow is documented in the Report on Use of AI Tools \cite{aiapp}.

\subsection{Assumptions}

\begin{enumerate}[label=\textbf{A\arabic*.},leftmargin=2.4em]
\item \emph{Stationary weekday regime.} The seven weekdays between 10 and 19 May
  that are fully covered by the sample are exchangeable draws from one daily regime.
  Weather, holidays and price promotions are not modelled. The observed daily volume
  varies with a coefficient of variation of 8.25\% around its mean (\figref{fig:data}a),
  so any single-day quantity in this paper carries that order of sampling noise.

\item \emph{Poisson trip generation.} Departures from cell $i$ in time interval
  $t$ are Poisson with rate $\lambda_{it}$, constant on 15-minute intervals. This is
  the standard assumption for individual trip-making; pooled across cells the counts
  are over-dispersed relative to Poisson, which reflects heterogeneity in the rates
  rather than a failure of the conditional model.

\item \emph{Supply-independent destination choice.} A rider chooses a destination
  without regard to whether a bicycle will be available there. This is the usual
  assumption of trip-distribution modelling. It is conditional here: the 50-zone LP
  finds no zone stockout at the measured inventory, but the cell-level simulation
  finds 5.84\% unserved at 343,638 and 3.51\% at 495,540, so the censoring is small
  but not zero and is reported rather than assumed away.

\item \emph{Endogenous attraction.} The attraction of a destination is measured by
  the arrivals it actually receives. The model is therefore a description of an
  equilibrium, not a structural model of land use: it may be used to predict how
  flows respond to \emph{distance} and to \emph{time of day}, but not to predict the
  consequence of changing what is built at a destination. We flag this again in
  Section~\ref{sec:eval}.

\item \emph{Homogeneous, reliable bicycles.} All bicycles in service are usable; a
  bicycle identifier persists until the bicycle is retired. The dataset reports a
  bicycle type, which is constant per bicycle and therefore carries no information
  for our purposes.

\item \emph{Truck operations.} A truck carries at most $Q=40$ bicycles, travels one
  inter-zone leg per hour, and unloads entirely at its destination. Loading and
  unloading are assumed to fit inside the same hour. Trucks may serve only the eight
  nearest zones within 9 km, a restriction that keeps the arc set to 290 and matches
  the range of an overnight shift.

\item \emph{Utility of accepting an incentive.} A rider offered $r$ CNY to end at a
  different cell evaluates $U = \kappa r - \rho\,\tau(d) + \epsilon$, with $\tau(d)$
  the extra travel and walking time for a detour $d$, and $\epsilon$ logistic. The
  parameters $\kappa,\rho$ and the inertia constant $c_0$ are not identified from our
  data: the order file contains no reward records, although Mobike had been operating
  red-envelope (hongbao) bikes since March 2017 \cite{technode2017,sixthtone2017}.
  They are therefore treated as behavioural parameters and swept in
  Section~\ref{sec:sens}.

\item \emph{Parking demand equals arrivals, on a standing stock.} A parked bicycle
  occupies space where its trip ended. Parking demand at a location is therefore the
  arrival process, on top of whatever bicycles already stand there; capacity is the
  peak of that path, not the peak of arrivals from an empty start.

\item \emph{Cost parameters.} A bicycle costs 1.0 CNY per day to own and maintain
  (about 1{,}000 CNY over a three-year life); a truck costs 600 CNY per day to own,
  insure and staff, and 6 CNY per kilometre to run. These are planning assumptions
  of the right order for 2017 Beijing. The truck-kilometre price was swept from 3 to
  12 CNY in an earlier run on the observed-OD surface; that run predates the T007
  corrections and was not repeated, so no price-sensitivity numbers are quoted. All
  costs are linear in the other parameters.
\end{enumerate}

\subsection{Notation}

\tabref{tab:notation} lists the symbols that recur. Two index sets run in parallel ---
cells $i,j$ for behaviour, parking and incentives, and zones $z$ for rebalancing ---
because the two models operate at different resolutions; where a quantity exists at
both, the subscript makes the resolution explicit. Time is indexed by the hour $t$ in
the optimisation and grouped into five bands $h$ where the decay is estimated, since
the behaviour model needs a rate at 15-minute resolution but a distance decay only at
the resolution of the daily rhythm.

\begin{table}[htbp]
\centering
\caption[Principal symbols]{Principal symbols. Indices are suppressed where the meaning is clear.}
\label{tab:notation}
\small
\begin{tabular}{@{}ll@{\hspace{1.2cm}}ll@{}}
\toprule
\textbf{Symbol} & \textbf{Meaning} & \textbf{Symbol} & \textbf{Meaning} \\
\midrule
$i,j$ & origin, destination cell & $\beta_h$ & distance decay in time band $h$ \\
$z$ & traffic zone ($Z=50$) & $d_{ij}$ & great-circle distance (km) \\
$t$ & hour ($T=24$) & $A_j$ & attraction of destination $j$ \\
$h$ & time band (five per day) & $P(j\mid i,h)$ & destination probability \\
$\lambda_{it}$ & departure rate (trips/h) & $n_{ij}$ & observed OD count \\
$p$ & order sampling rate & $\mu$ & rides per bicycle per day \\
$F$ & fleet size (bicycles) & $D$ & distinct bicycles observed \\
$I_{zt}$ & inventory of zone $z$ at $t$ & $u_{zt}$ & unmet departures \\
$x_{ijt}$ & bicycles moved $i\to j$ in $t$ & $m_{ijt}$ & trucks moved $i\to j$ in $t$ \\
$K$ & trucks in service & $\alpha$ & fleet size relative to observed \\
$q(r)$ & placements bought by reward $r$ & $r$ & reward (CNY) \\
$R$ & catchment radius (km) & $w_c$ & parking capacity (slots) \\
\bottomrule
\end{tabular}
\end{table}
